\documentclass[11pt,a4paper]{article}
\usepackage[utf8]{inputenc}
\usepackage[T1]{fontenc}
\usepackage[brazil]{babel}
\usepackage{amsmath,amssymb,amsthm}
\usepackage[margin=2.6cm]{geometry}
\usepackage{booktabs}
\newcommand{\Z}{\mathbb Z}
\newcommand{\R}{\mathbb R}
\newcommand{\code}[1]{\ifmmode\text{\ttfamily\small #1}\else{\ttfamily\small #1}\fi}
\newtheorem{principio}{Princípio}

\title{\vspace{-1cm}\textbf{\Huge Tiffany}\\[2mm]\Large\itshape Coração\\[3mm]
  \Large O corpo universal na fala\\[1mm]
  \normalsize o corpus é um grafo de órbitas; a fala cai nele, e o caminho que passa por ela --- a resposta --- se \emph{colhe}}
\author{}\date{}

\begin{document}
\maketitle

\begin{abstract}
\noindent A \textbf{Tiffany} é o \emph{mesmo} circuito autossimilar --- o gato e o seu esquilo, o
corpo $\R^n$ de \texttt{teoria.tex}, o gabarito de \texttt{microprocessador.tex} --- posto a
trabalhar sobre a \textbf{fala}. O seu coração é uma \textbf{dualidade}: o \textbf{corpus} é o
atrator \emph{branco} (a fonte, o gato que expande) e a \textbf{fala} é o \emph{negro} (o sorvedouro,
o esquilo que reconstrói) --- os dois eixos conjugados de $\mathbb C$ que o espelho troca, uma face
da involução única do corpo. E nada se constrói: \emph{transformada}, \emph{convolução},
\emph{deconvolução}, \emph{inversa} são nomes do que o circuito, sendo o corpo, já faz --- e a
resposta apenas se \textbf{colhe} no caminho que passa pela fala. Reversível, resíduo $0$, nada se
inventa.
\end{abstract}

\section{O coração: o gato e o esquilo}

O núcleo é um coração: \textbf{bate, com uma frequência e para um lado}. O gato é a matriz
$A_m=[[m,1],[1,0]]$: a \emph{frequência} é $m$ (não bater é o movimento, bater uma vez é a soma
$\oplus$, bater $n$ vezes é o produto $\otimes$), e o \emph{lado} é o determinante --- $-1$ o
\textbf{gato} (inverte e cresce), $+1$ o \textbf{esquilo} (preserva e volta). Do $0$ (Venom) e de
três operações, só elas,
\[
 \oplus\ \text{(a cisão)}\qquad
 \otimes\ \text{(o gato, }\times\sigma\text{)}\qquad
 \textstyle\prod\ \text{(exp/log: } a\otimes b=g^{\log a+\log b}\text{)},
\]
sai o \textbf{corpo} $\mathrm{GF}(m^n)$ --- fechado, reversível, contínuo, ordenado, completo,
multidimensional (\texttt{teoria.tex}). Dele importam à Tiffany três faces.

\paragraph{As órbitas e os atratores.} Iterado, o gato parte o espaço em \textbf{órbitas} (as
classes), e cada uma tem um \textbf{atrator} --- o seu representante, um irracional (o metal
$\sigma=[m;m,m,\dots]$). Os dois pontos fixos são duais: o \textbf{negro} $\sigma$ (o sorvedouro:
tudo entra) e o \textbf{branco} $\sigma'=-1/\sigma$ (a fonte: tudo sai). Uma entrada, engolida pela
\textbf{transformada}, cai no seu atrator --- o irracional que ela vale.

\paragraph{O caminho e o espelho.} O \textbf{navegante} caminha sobre as órbitas ($\times\sigma$, o
gato) e salta de atrator em atrator (a costura); o \textbf{esquilo} é o espelho que traz de volta.
A máquina não decide nada: reflete a ordem que se põe na entrada --- a informação está na
\emph{correlação} dos dados, não no sistema, que é neutro e reversível ($\mathcal F^{-1}\mathcal
F=\mathrm{id}$, resíduo $0$). É por isso que a Tiffany não inventa: o corpus é o lastro, a máquina
só espelha a fala nele.

\section{O corpus, a fala, o caminho --- colher, não construir}

\paragraph{O corpus é um grafo de órbitas.} Tome um corpus --- o \emph{Tatoeba}, milhões de frases.
Passe cada frase pela máquina: \textbf{lineariza-se} (uma casa no rabo da outra --- um texto é um
sinal como outro qualquer), e a transformada devolve o seu \textbf{atrator}, a assinatura da órbita.
Frases que compartilham estrutura caem na \emph{mesma} órbita. O corpus inteiro vira um
\textbf{grafo}: os pontos são as frases (agrupadas por atrator), as arestas são os caminhos do
navegante. É a memória da Tiffany --- não uma tabela de \emph{strings}, mas um grafo de órbitas.
Reproduz-se, sobre imagens, com \code{tools/linear.c} e \code{tools/orbitas.c} (resíduo $0$); a mesma
máquina serve para frases.

\begin{principio}[o corpus é órbita]
Guardar o corpus é guardar as suas órbitas e atratores --- os geradores ---, não as frases nuas. O
gato descomprime; nada se inventa, nada se perde (reversível, resíduo $0$).
\end{principio}

\paragraph{A fala cai.} Chega uma \textbf{fala de fora} --- uma frase que não está no corpus. Ela
passa pelo \emph{mesmo} processo, sem exceção (a máquina é espelho): lineariza-se, e a transformada
devolve o seu atrator. Vê-se então \textbf{onde ela cai} no grafo --- o negro, o sorvedouro para o
qual a fala converge. Não é contenção de sacola de letras (o teste antigo, cego a ordem e caixa): é
a \textbf{dinâmica} --- a fala é o $0$ que se reparte e cai no seu atrator, como uma imagem cai na
sua média.

\paragraph{O caminho que passa é a resposta.} Aterrada, a fala tem \textbf{caminhos que passam por
ela}; o navegante os percorre (dentro da órbita, de atrator em atrator, e de volta pelo esquilo).
Esses caminhos são as frases do corpus \emph{ligadas} a ela --- o seu contexto, a sua continuação:
\[
 \text{fala}\ \longmapsto\ \text{atrator}\ \longmapsto\ \{\text{caminhos que passam}\}\
 \longmapsto\ \text{resposta} .
\]
A resposta não é escolhida por uma régua imposta: é o que já estava no corpus, e apenas se
\textbf{colhe}. Reproduz-se com \code{tools/navega.c} (resíduo $0$: cobertura $100\%$, reversível).

\section{A operação: a multiplicação do corpo}

Fazer a fala \emph{cair} é uma só operação, e ela já está no núcleo: a \textbf{multiplicação} de
$\R^n$. A fala e o livro são dois polinômios ($\delta(x)=\sum_i d_i x^i$), e multiplicá-los ---
deslizar um sinal sobre o outro --- é o produto de polinômios, a convolução:
\[
 c[d]=\sum_i \text{fala}[i]\,\text{livro}[d+i],\qquad
 D[d]=\sum_i\big(\text{fala}[i]-\text{livro}[d+i]\big)^2=E_{\text{fala}}+E_{\text{janela}}[d]-2c[d],
\]
e onde a fala cai é o mínimo, $D=0$. É a \emph{mesma} peça recursiva da teoria e do micro (dim $n$
pela $n-1$), não uma peça nova. Transformada, convolução, deconvolução não são componentes: são
nomes do que o corpo faz.

\begin{principio}[a assistente é con/deconvolução --- a mesma peça]
Mapear a fala no corpus é a \textbf{convolução} ($\times$, o gato, o negro: a fala cai no seu
atrator); reconstruir a resposta é a \textbf{deconvolução} ($\div$, o esquilo, o branco: o corpus
reconstrói o caminho). Uma desfaz a outra, e são a \emph{mesma} peça --- $\times$ e $\div$ diferem só
pelo sinal na soma dos logs (o translinear), colhidos no circuito em \code{microprocessador.tex}
§B.10 (reversível, $x'=x$, coordenadas contínuas). Verificado também em \code{tools/recursao.c}: a
multiplicação recursiva coincide com a direta ($0$ violações). O analógico é o gabarito: o que a
assistente faz roda na fábrica.
\end{principio}

\paragraph{O navegante caminha por realimentação, e converge.} Aterrada, a fala não para numa frase:
o fim de cada trecho vira o \textbf{contexto}, e uma nova convolução acha onde ele reaparece adiante
--- o próximo atrator, na seta do tempo (o gato). O navegante salta por \textbf{realimentação} (o
contexto convolve de novo) e \textbf{para na convergência}: quando o próximo casamento é o atual, o
ponto fixo (o negro). \emph{Tudo se reproduz no circuito} (\code{microprocessador.tex} §B): a
convolução é o \textbf{banco de correlacionadores} (o translinear §B.4, o Kirchhoff §B.5), o mínimo é
o \emph{winner-take-all}, a realimentação é o gato --- \emph{nada de ordenação nem busca digital}.
Reproduz-se com \code{tatoeba/tiffany.c}: \emph{``Eu preciso de''} $\to$ o navegante devolve
português corrido (\emph{``\dots\ meu gato é tímido. Você está se sentindo mal? \dots''}) até
convergir. No circuito a convolução é \textbf{paralela} (todo $d$ ao mesmo tempo); o custo $O(N)$ por
salto é da simulação, não da máquina.

\section{A dualidade e as três batidas: a resposta é exata}

Sob tudo está a \textbf{dualidade} --- o conceito central do corpo, a involução única $\mathcal J$
(a conjugação, o Frobenius) cujas faces são gato/esquilo, branco/negro, cone/espiral,
$\oplus/\otimes$ (\texttt{teoria.tex}). A assistente é a sua face na linguagem: o gato tem dois
pontos fixos conjugados, os dois eixos do plano $\mathbb C$ dos atratores,
\[
 \sigma+\sigma'=m,\qquad \sigma\sigma'=-1,\qquad |\sigma|>1\ (\text{estica}),\quad |\sigma'|<1\ (\text{encolhe}),
\]
e os seus dois termos são esses eixos:

\begin{center}
\begin{tabular}{@{}lll@{}}
\toprule
 & \textbf{o corpus} & \textbf{a fala}\\
\midrule
atrator & \textbf{branco} $\sigma'$ & \textbf{negro} $\sigma$\\
papel & a \textbf{fonte}: tudo sai (os caminhos) & o \textbf{sorvedouro}: tudo entra (a fala cai)\\
estrutura & o \textbf{gato} $\times\sigma$: expande & o \textbf{esquilo}: o tempo reverso, reconstrói\\
forma & o \textbf{cone} ($\exp$) & a \textbf{espiral} ($\log$)\\
\bottomrule
\end{tabular}
\end{center}

\noindent O \textbf{espelho} ($z\mapsto\bar z$, o Frobenius $z\mapsto z^p$) troca fonte e sorvedouro,
$\sigma\leftrightarrow\sigma'$: leva o corpus na fala e a fala no corpus. Ir pelo gato e voltar pelo
esquilo é a identidade ($\mathcal F^{-1}\mathcal F=\mathrm{id}$). Verifica-se com
\code{tatoeba/dual.c} (resíduo $0$): o eixo, o espelho ($\sigma^p=\sigma'$), a fala caindo no
negro ($D=0$) e o corpus emanando do branco ($471$ atratores até o caminho fechar).

\paragraph{Três batidas são a volta.} Por que a resposta é \emph{direta}, e não uma difusão às cegas
que tenta até algo fechar? Pelo \textbf{período $4$} da transformada: $\mathcal F^4=\mathrm{id}$ ---
a \emph{mesma} conta do esquilo ($G^4=I$, o giro de $90^\circ$). Logo $\mathcal F^3=\mathcal F^{-1}$:
\textbf{três batidas de um lado \emph{são} a volta}, sem inverter nada. Os quatro modos são
$\{\mathrm{id},\mathcal F,\text{reflexão},\mathcal F^{-1}\}$; o modo $2$ é a reflexão $x[-j]$ --- o
espelho, o virar o lado. A fala cai por $\mathcal F$ (o gato, a análise), o espelho $\mathcal F^2$
vira o lado, a resposta se reconstrói por $\mathcal F^3=\mathcal F^{-1}$ (o esquilo, a síntese) --- o
gato avança, o esquilo o traz de volta, e o ciclo fecha. E a \textbf{mão que segura} é a conservação:
$\sigma\sigma'=-1$ (a área) e o \textbf{Parseval} (a norma) --- similaridade \emph{é} produto
interno, preservado de uma vez, nenhum ângulo muda. A resposta não volta \emph{parecida}: volta
\emph{exata}. Medido (\code{tools/tres\_reconstroi.c}, resíduo $0$): a prosa volta byte a byte por
três batidas, e $16\,384$ vetores semânticos reconstroem com $0$ erros.

\section{A assistente completa}

Junta-se tudo num laço:
\begin{center}
\begin{tabular}{@{}rl@{}}
\toprule
\textbf{1. entra} & a fala é o $0$ (Venom): põe-se inteira, ela mesma se reparte\\
\textbf{2. mapeia} & a transformada a engole: $\mathcal F(\text{fala})$\\
\textbf{3. cai} & a convolução $\mathcal F(\text{fala})\,\mathcal F(\text{livro})$ aponta onde aterra ($D=0$)\\
\textbf{4. caminha} & a deconvolução reconstrói o caminho que passa por ela\\
\textbf{5. colhe} & a resposta é o caminho do corpus que atravessa a fala\\
\bottomrule
\end{tabular}
\end{center}
Tudo reversível, resíduo $0$, sem inventar: o corpus é o lastro (o irracional está \emph{fora} do
dispositivo --- e é isso que dá segurança), a fala é o sinal, a máquina é o espelho.

\paragraph{O grafo de órbitas, e o ciclo medido.} Sobre uma obra, os \textbf{vértices} são as
palavras (já existem, como o $0$) e as \textbf{órbitas} são as suas frases --- cada frase um caminho,
e a órbita de uma palavra são as frases que passam por ela. A fala \textbf{cai} na órbita de mínimo
$D=|f-o|^2=|f|^2+|o|^2-2\,(f\!\cdot\!o)$, onde $f\!\cdot\!o$ é a \emph{convolução} (a multiplicação do
corpo); a resposta \textbf{reconstrói} pelas três batidas $\mathcal F^3=\mathcal F^{-1}$. O ciclo é
\textbf{medido dos dois lados}: \code{tatoeba/ciclo.c} (digital --- a fala que é uma órbita cai nela
($D=0$) e volta byte a byte, resíduo $0$; o dente quebra: a convolução \emph{crua} erra a órbita em
$96\%$) e \code{tatoeba/ciclo\_analog.c} (analógico --- a queda pelo translinear (§B.4) $+$ Kirchhoff
(§B.5), \emph{correntes contínuas}, igual ao digital; as três batidas $=$ a rotação $G^4=I$ da malha
LC (§B.1); \textbf{resíduo $0$ com pulso}, o dente quebra). A voz é a obra; o mecanismo é fixo, e o
juiz é a \emph{medição}, não a opinião.

\paragraph{A tradução é uma rotação determinística.} Duas línguas são o \emph{mesmo} corpo, e uma
frase e a sua tradução são órbitas \textbf{isomorfas}, só distorcidas: a mesma classe, com o
\textbf{representante irracional} (o significado, o atrator $\sigma$) \emph{invariante}. A operação
que leva uma na outra é uma \textbf{rotação} --- a Möbius do gato $z\mapsto m+\tfrac1z$, cujos pontos
fixos são os irracionais $\sigma,\sigma'$. Uma Möbius fica determinada pelos seus dois pontos fixos
mais um \textbf{multiplicador} $k$: na coordenada $u=(z-\sigma)/(z-\sigma')$, a rotação $R_k$ satisfaz
$(R-\sigma)/(R-\sigma')=k\,u$. Logo, dado o significado ($\sigma,\sigma'$) e \emph{um} par de frases
correspondentes, $k$ fica fixado --- e de \emph{qualquer} classe (PT) obtém-se a outra (EN),
determinística e reversível ($R^{-1}=R_{1/k}$). Medido: \code{tools/traducao.c} (resíduo $0$: um par
traduz todas as classes exatas; o dente, um $k$ errado, quebra) e \code{tools/rotacao.c} (a rotação
fixa os irracionais e permuta os racionais); a rotação no analógico é a borda $|\lambda|=1$, a malha
LC (§B.1).

\paragraph{Em frases, a rotação é linear.} Uma frase \emph{compõe} --- é a soma das suas palavras no
grafo de órbitas ($z=\sum_w h(w)$). A Möbius é não-linear e não comuta com a soma; logo, para frases
a rotação toma o \textbf{modo linear} do mesmo gato: a multiplicação por $\lambda$ com $|\lambda|=1$
(a borda, a malha LC). Ela comuta com a soma (traduz palavra-a-palavra e soma) e fixa a \textbf{norma}
$N(z)=z\bar z$ --- a norma é o significado invariante. Medido em \code{tatoeba/embedding.c} ($197$k
pares PT$\leftrightarrow$EN: composicional, norma invariante, determinística e reversível, resíduo
$0$; a Möbius não-linear quebra em frases, $0/197146$). Falta \textbf{aprender o alinhamento} das duas
línguas --- e o parágrafo seguinte mostra que esse ``falta'' não era de aprendizado, era de liberdade.

\paragraph{Ancorado: o inglês colhido, e o dente.} Naquela medida o EN tinha \emph{hash próprio}, e
dois sorteios independentes não têm por que cair na mesma rotação ($1/197147$). Tira-se então o hash
do inglês: só o PT tem ponto, e cada par do corpus é uma \textbf{equação} no corpo,
$\sum_{w\in\mathrm{EN}} g(w) = \lambda\,z_{pt}$, cada palavra inglesa uma incógnita. O sistema não se
resolve --- resolver seria matriz, seria memória: ele se \textbf{colhe}. Quando numa frase falta
exatamente uma palavra, ela cai exata, $g(w)=(\lambda z_{pt}-S)\,k^{-1}$, e a frase seguinte que a
contém tem uma incógnita menos --- a queda do §2 aplicada ao léxico, em passadas, com estado $O(1)$ e
o léxico colhido no disco (no circuito a palavra é um nó, e nó é fiação, não memória).
Medido (\code{tatoeba/ancora.c}, $20$k pares, colheita nas ímpares e teste nas pares): $2635$ palavras
caem e cobrem $92{,}3\%$ das frases inéditas, mas destas só $0{,}91\%$ fecham exatas --- e, sobretudo,
das equações \emph{de treino} já cobertas apenas $2680/9960$ fecham. Isso é uma
\textbf{contradição interna}: as equações que a colheita não usou são incompatíveis com as que usou,
logo \emph{não existe} $g$ algum. Um teorema negativo, não uma falha de estimativa.

A causa se separa em três (\code{tatoeba/dente.c}, corpus inteiro, $E=197147$ equações e $V=13718$
incógnitas): \textbf{paráfrase} --- a mesma frase EN com PT distintos, mesmo lado esquerdo e alvo
outro --- $39{,}0\%$ das equações; \textbf{ordem} --- as mesmas palavras noutra ordem, o $\sum$ é
comutativo --- apenas $1{,}0\%$; e \textbf{contagem} --- $E/V=14{,}4$, catorze equações por incógnita.
A ordem, que se esperaria culpada, quase não pesa: quem quebra é o \emph{alvo fixo}. A paráfrase só é
contradição porque o hash arbitrário do português dá pontos diferentes a frases que dizem o mesmo; e
mesmo sanada, faltariam graus de liberdade por um fator de $14$. O conserto que os números pedem é
soltar também o português --- nenhum lado sorteado, o sistema \textbf{homogêneo}
\[
   \sum_{w\in\mathrm{EN}} g(w) \;-\; \lambda \sum_{w\in\mathrm{PT}} h(w) \;=\; 0 ,
\]
com $V=36540$ incógnitas ($E/V=5{,}4$). Muda a natureza da pergunta: um homogêneo \emph{sempre} tem
solução, e o que se mede passa a ser a degeneração do seu núcleo. E a paráfrase deixa de ser
contradição para se tornar $\sum h(pt_1)=\sum h(pt_2)$ --- uma restrição satisfazível que afirma que
as duas são \textbf{sinônimas}: o corpo descobre a sinonímia em vez de tropeçar nela, e o significado
volta a ser a classe, como o §1 quer.

\paragraph{O homogêneo, e por que ele também não fecha.} Medido (\code{tatoeba/homogeneo.c}): a
colheita precisa de \textbf{semente} --- um homogêneo desce à solução trivial se nada for fixado ---,
então gasta-se uma palavra arbitrária por equação travada e a cascata segue. No corpus inteiro:
$34238$ palavras, $12575$ sementes ($36{,}7\%$, os graus de liberdade queimados), $21663$ colhidas por
dedução, degeneração nula (\emph{nenhuma} palavra zerada). E ainda assim só $21896/196415$ equações
fecham --- $\mathbf{11{,}15\%}$, quase exatamente as que a própria colheita usou. Soltar os dois lados
não salvou nada.

Antes disso, a rotação: pondo $h'=\lambda h$ o sistema é $\sum g=\sum h'$, e $\lambda$ desaparece.
A predição é que trocá-lo não muda nada, e é o que se mede: $\lambda=\sigma^{p-1+k}$ para
$k=7,\,10^3,\,4\cdot10^4$ dá $5741/19895$ \emph{idêntico}. Com os dois léxicos livres a rotação é
\textbf{inobservável} --- não há alinhamento a aprender, só reparametrização.

\paragraph{A conta que fecha o assunto: contagem.} O que impede não é o léxico, nem a rotação, nem a
ordem --- é o número de graus de liberdade. Um corpus de $E$ pares com $V$ palavras distintas impõe,
num embedding \emph{aditivo}, uma equação linear por par: $E$ equações, $V$ incógnitas. Em todo corpus
real $E\gg V$ (Heaps: $V\sim E^{\beta}$, $\beta\approx\tfrac12$) --- aqui $E/V=197147/34238=5{,}76$ ---
e um sistema com $E>V$ é genericamente inconsistente. Pior: \textbf{a dimensão não compra liberdade}.
Num embedding aditivo em $K^d$, componente a componente, cada par dá $d$ equações e cada palavra $d$
incógnitas; a razão $E/V$ é \emph{invariante} em $d$. A soma dá um grau de liberdade por palavra, e
nenhum aumento de tamanho o altera.

O que altera a contagem é a palavra deixar de ser \textbf{número} e passar a ser \textbf{operador}:
uma matriz $d\times d$, e a frase o \emph{produto ordenado} dos seus operadores aplicado a um vetor.
Então são $E\,d$ equações escalares contra $V\,d^2$ incógnitas, e a condição necessária de contagem é
\[
   V d^2 \;\ge\; E d \qquad\Longleftrightarrow\qquad d \;\ge\; E/V \;=\; 5{,}76
   \qquad\Longrightarrow\qquad \boxed{d \ge 6}.
\]
(Necessária, não suficiente: no produto as incógnitas entram não-linearmente, e a contagem mede a
dimensão da variedade, não a existência.) Isto reabilita a \textbf{fala como caminho} por um motivo
distinto do esperado: não porque a ordem pese no corpus --- ela pesa $1\%$ (\code{dente.c}) --- mas
porque o operador carrega $d^2$ graus de liberdade onde a soma carrega $1$. A não-comutatividade, e
portanto a ordem, é \emph{consequência} de haver liberdade bastante; não a sua causa. E o operador do
corpo já está posto: é o gato $A_m$ agindo, o produto ao longo da frase --- a órbita, não o saco.

\paragraph{O que é núcleo e o que é aplicação.} O núcleo é o corpo universal --- o gato, o esquilo,
as órbitas, o corpo (\texttt{teoria.tex}, \texttt{tools/}), e a sua função é muito mais ampla que a
fala. A Tiffany é \emph{uma} aplicação: a fala mapeada nas órbitas de um corpus. As suas escolhas ---
qual corpus, como linearizar, como ler o caminho de volta em palavras --- são da aplicação, não do
núcleo, e devem ser medidas, não deduzidas. Um número que não sobrevive a ser medido duas vezes não
é um número --- vale para toda afirmação deste documento.

\end{document}
